%% BioXRep --- arXiv short paper (ACL / BioNLP style, non-anonymous)
%% Cross-Notation Biomedical Entity Normalization
%%
%% Build:  pdflatex main -> bibtex main -> pdflatex main -> pdflatex main
%% Requires the ACL style files (acl.sty, acl_natbib.sty, acl_natbib.bst) from
%%   https://github.com/acl-org/acl-style-files
%% or Overleaf's official "ACL ARR / *ACL Paper Template" gallery entry.
%% This is the arXiv preprint form: authors are shown (no [review] option).
%% For a BioNLP/ACL 2027 anonymous submission, change to \usepackage[review]{acl}.
\documentclass[11pt]{article}

\usepackage[]{acl}
\usepackage{times}
\usepackage{latexsym}
\usepackage[T1]{fontenc}
\usepackage[utf8]{inputenc}
\usepackage{microtype}
\usepackage{inconsolata}
\usepackage{booktabs}
\usepackage{graphicx}
\usepackage{multirow}

\title{Cross-Notation Biomedical Entity Normalization:\\
       A Contrastive Char-CNN Student and What It Fails to Learn}

\author{First Author \\
  Affiliation \\
  \texttt{first@domain} \\\And
  Second Author \\
  Affiliation \\
  \texttt{second@domain} \\}

\begin{document}
\maketitle

\begin{abstract}
Biomedical entities appear under many surface notations---approved gene symbols,
alias symbols, prior symbols, and free-text descriptive names---and mapping among
them is a core normalization problem. We ask whether a lightweight character-level
convolutional encoder, trained with a supervised contrastive loss over gene
equivalence classes, can learn a cross-notation representation that transfers to a
notation held out of training. On an HGNC alias-retrieval benchmark with a
fact- and notation-disjoint split, the student reaches $0.985$ in-domain
validation top-1 but only $0.059 \pm 0.001$ MRR (mean $\pm$ std over $5$ seeds) on
held-out \texttt{alias\_symbol} queries
against a $45$k-way pool---below a character $n$-gram lexical floor ($0.077$) and a
SapBERT dense retriever ($0.134$). We report this negative result alongside four
reference methods (character $n$-gram, SapBERT dense, a BioSyn-style sparse--dense
hybrid, and a structured canonical-teacher oracle) evaluated on an identical pool.
The analysis is consistent across three query notations: an off-the-shelf
biomedical embedding wins on low-overlap descriptive names, lexical and hybrid
methods win on high-overlap prior symbols, and a small surface-only student beats
none of them on the disjoint task it was built to contest. An
attention-distillation ablation shows that supervising the student with a stronger
surface-form teacher's attention---deterministic or neural---does not close the gap
either. We argue the fact-disjoint construction is what exposes the gap, and we
outline what closing it requires.
\end{abstract}

\section{Introduction}
\label{sec:intro}
Integrating biomedical information requires recognizing that the same underlying
object is written many ways. A single gene is referred to by an approved symbol,
one or more alias symbols, former official symbols, and free-text descriptive
names; a single variant is written in protein and nucleotide HGVS, in database
identifiers, and in prose. Normalizing these surface notations onto a shared
representation is a prerequisite for linking mentions across resources, and neural
encoders are known to be sensitive to exactly the surface details---numeric and
notational---that normalization must abstract away
\citep{wallace2019, xval, thawani2021representing}.

A natural hypothesis is that distilling a compact student against a canonical
target---the transfer-of-knowledge paradigm that trains a small model to match a
richer supervisory signal \citep{bucila2006, hinton2015distilling}---can teach the
student to collapse equivalent notations onto one point in embedding space. We test
the strongest lightweight version of this hypothesis: a
byte-level character CNN ($\sim$100k parameters) trained with a supervised
contrastive loss \citep{supcon} that pulls same-gene forms together, using the
HGNC approved symbol as the canonical anchor. We then evaluate it under a
deliberately hard protocol in which an entire query \emph{notation} is held out of
training, so success cannot come from having seen the target surface form.

Our contribution is a controlled negative result. The student trains correctly
(in-domain validation top-1 $0.985$) yet fails to transfer: it scores
$0.059 \pm 0.001$ MRR (mean $\pm$ std over $5$ seeds)
on the held-out \texttt{alias\_symbol} task, below both a character $n$-gram floor
and an off-the-shelf SapBERT dense retriever. We situate this result among four
reference methods on an identical candidate pool and across three query notations,
and we show that the fact- and notation-disjoint split is what makes the
generalization gap visible where an in-distribution measurement would hide it.

\section{Related Work}
\label{sec:related}
\paragraph{Biomedical entity normalization.} Synonym-aware and self-aligned
encoders set the current standard for linking surface mentions to canonical
concepts. BioSyn \citep{biosyn} learns to combine sparse and dense similarity for
biomedical entity linking, and SapBERT \citep{sapbert} self-aligns a biomedical
language model over UMLS synonym pairs. We use both as off-the-shelf learned
baselines.
\paragraph{Variant nomenclature.} Text mining and standardized representation of
genetic variants---tmVar and its successors \citep{tmvar, tmvar2, tmvar3} and the
GA4GH Variation Representation Specification \citep{ga4ghvrs}---frame variant
equivalence as canonicalization across free text and structured identifiers,
the same cross-notation problem in a different modality.
\paragraph{Contrastive representation learning.} Instance-level and class-level
contrastive objectives \citep{simclr, supcon} and language-supervised alignment of
paired views \citep{clip} underpin the student's training signal: distinct
notations of one gene are treated as views to be aligned.
\paragraph{Knowledge distillation.} Training a compact student to reproduce the
behavior of a stronger model dates to model compression \citep{bucila2006} and its
soft-label formulation \citep{hinton2015distilling}, and has been extended from
matching outputs to matching intermediate representations
\citep{ba2014deep, romero2015fitnets}, attention maps \citep{zagoruyko2017attention},
and structural relations via contrastive objectives \citep{tian2020crd}; see
\citet{gou2021survey} for a survey. In NLP the paradigm compresses large
pretrained encoders into task-agnostic students \citep{sanh2019distilbert,
jiao2020tinybert, wang2020minilm}, and in dense retrieval it transfers relevance
signal across architectures \citep{hofstatter2020improving} and languages
\citep{reimers2020multilingual}. Our student is trained with a supervised
contrastive loss anchored on a canonical symbol; we additionally test an
attention-distillation objective \citep{zagoruyko2017attention} against both a
deterministic and a trained neural teacher (\S\ref{sec:ablation}), and find neither
closes the cross-notation gap. We treat these methods as the reference for what a
\emph{semantically pretrained} teacher-strength signal---which our ablation does not
supply---would add.
\paragraph{Numeric and symbolic representation.} A parallel line studies how models
represent the numeric content of text. Subword tokenizers fragment numbers
inconsistently, and probing shows encoders capture magnitude only weakly
\citep{wallace2019, thawani2021representing}. Proposed remedies include
numeracy-preserving word embeddings \citep{spithourakis2018numeracy,
sundararaman2020numeracy}, explicit arithmetic units \citep{trask2018nalu},
injected numerical-reasoning skills \citep{geva2020injecting}, and continuous or
frequency-based value encodings---sinusoidal position features
\citep{vaswani2017attention}, random Fourier features \citep{tancik2020fourier}, and
continuous numeric tokenization \citep{xval}. Gene-symbol normalization is the
non-numeric half of the cross-notation problem; the numeric half---HGVS positions,
laboratory values, unit conversions---is where these value-aware encoders become the
natural tools, and we treat them as the reference point for extending the benchmark
beyond symbols. On an HGVS cross-notation retrieval task with position-confounded
hard negatives, we find that none of these value-aware input encodings---an explicit
normalized scalar, sinusoidal features, or xVal-style continuous tokenization
\citep{xval}---improves hard top-1 over a text-only encoder, and neither does an
auxiliary numeric-consistency loss. On the full candidate pool the numeric features
appear to help, but that gain is an artifact of trivially rejecting random-fill
decoys; when the positive is scored against only its position-matched hard negatives
the ranking inverts and text-only is strongest, because the shared position is a
confound rather than a discriminator. The numeric half thus inherits the same
capacity-limited negative as the symbolic half.
\paragraph{Clinical sequence models.} Representation learning over coded EHR
sequences \citep{medbert, behrt, clmbr, gbert} shares the multi-notation setting
but targets temporal streams of clinical codes rather than surface-form
normalization of entity strings.

\section{Task and Data}
\label{sec:data}
We build gene-alias equivalence classes from HGNC, where each class is one gene and
its members are the gene's forms under several notations. Relevance is defined by a
shared gene identifier (\texttt{fact\_id}). The retrieval task is: given an
alias-form query, retrieve the gene's \texttt{approved\_symbol} from a shared
candidate pool. We evaluate three query notations against one pool:

\begin{itemize}
\itemsep0.2em
\item \texttt{alias\_symbol} (held-out test split): the fact- and notation-disjoint
generalization task.
\item \texttt{prev\_symbol}: former official symbols, with high surface overlap with
the approved symbol.
\item \texttt{alias\_name}: free-text descriptive names, with low surface overlap.
\end{itemize}

The candidate pool is the $44{,}997$ approved symbols from the training split.
Queries are capped at $2{,}000$ per task, drawn from $44{,}718$ available held-out
\texttt{alias\_symbol} rows. Splits are fact-disjoint: the gene identifiers used to
form the \texttt{alias\_symbol} test set are excluded from training, so no test
gene is seen in any notation during training. All methods are scored on the
identical $($queries, pool$)$ pair per task, making within-task comparison direct.

\section{Method}
\label{sec:method}
The student is a byte-level character CNN (hidden dimension $64$, projection
dimension $128$, $\sim$100k parameters) that maps a surface string to an
$L_2$-normalized embedding. It is trained for $3$ epochs with a supervised
contrastive loss \citep{supcon} (temperature $0.07$) using class-aware batches over
$22{,}364$ HGNC equivalence classes, with the \texttt{alias\_symbol} test genes
excluded. Critically, the \texttt{alias\_symbol} notation is \emph{entirely absent}
from training: the student only ever sees \texttt{approved\_symbol},
\texttt{prev\_symbol}, \texttt{alias\_name}, \texttt{prev\_name}, and
\texttt{approved\_name} forms. Retrieval scores queries against the pool by cosine
similarity of these embeddings, through the same flat-pool evaluation harness used
for every baseline. Training loss decreases monotonically ($2.69 \rightarrow 1.41$)
and in-domain symbol$\leftrightarrow$name validation top-1 reaches $0.985$,
confirming the objective is learned.

\section{Baselines}
\label{sec:baselines}
We compare against four reference methods on the identical pool:
\textbf{character $n$-gram}, the cosine of character $3$-gram TF--IDF vectors, a
purely lexical floor; \textbf{SapBERT dense} \citep{sapbert}, the cosine of
$L_2$-normalized CLS embeddings; \textbf{BioSyn-style hybrid} \citep{biosyn},
$(1-\lambda)\cdot\text{sparse} + \lambda\cdot\text{dense}$ with $\lambda=0.7$; and
the \textbf{canonical teacher}, an exact match on the structured HGNC
\texttt{symbol} attribute that scores $1.000$ by construction and serves as a
structured ceiling rather than a learned competitor.

\section{Results}
\label{sec:results}
Table~\ref{tab:results} reports top-1, top-5, and MRR for all methods on the three
query notations. The student's only fact- and notation-disjoint measurement is
\texttt{alias\_symbol}; its \texttt{prev\_symbol} and \texttt{alias\_name} rows are
in-distribution (train-seen) and are marked accordingly.

\begin{table}[t]
\centering
\small
\setlength{\tabcolsep}{4pt}
\begin{tabular}{llccc}
\toprule
\textbf{Notation} & \textbf{Method} & \textbf{top-1} & \textbf{top-5} & \textbf{MRR} \\
\midrule
\multirow{5}{*}{\shortstack[l]{\texttt{alias\_}\\\texttt{symbol}\\(held-out)}}
 & Char $n$-gram        & 0.047 & 0.108 & 0.077 \\
 & SapBERT dense        & 0.080 & 0.190 & \textbf{0.134} \\
 & BioSyn hybrid        & 0.077 & 0.182 & 0.129 \\
 & Char-CNN student$^{\ddagger}$ & 0.041 & 0.075 & 0.059 \\
 & Canonical teacher    & 1.000 & 1.000 & 1.000 \\
\midrule
\multirow{5}{*}{\texttt{prev\_symbol}}
 & Char $n$-gram        & 0.146 & 0.276 & 0.205 \\
 & SapBERT dense        & 0.140 & 0.297 & 0.219 \\
 & BioSyn hybrid        & 0.175 & 0.325 & \textbf{0.250} \\
 & Char-CNN student$^{\dagger}$ & 0.118 & 0.203 & 0.158 \\
 & Canonical teacher    & 1.000 & 1.000 & 1.000 \\
\midrule
\multirow{5}{*}{\texttt{alias\_name}}
 & Char $n$-gram        & 0.026 & 0.056 & 0.042 \\
 & SapBERT dense        & 0.234 & 0.404 & \textbf{0.316} \\
 & BioSyn hybrid        & 0.225 & 0.398 & 0.308 \\
 & Char-CNN student$^{\dagger}$ & 0.026 & 0.053 & 0.042 \\
 & Canonical teacher    & 1.000 & 1.000 & 1.000 \\
\bottomrule
\end{tabular}
\caption{Retrieval of the \texttt{approved\_symbol} over a shared $44{,}997$-symbol
pool ($2{,}000$ queries per notation). $^{\dagger}$train-seen: the student saw this
query notation during training, so these are in-distribution measurements, not
held-out generalization. Best learned MRR per notation in \textbf{bold}; the
canonical teacher is a structured oracle upper bound.
$^{\ddagger}$the held-out \texttt{alias\_symbol} student row is the mean over $5$
seeds ($\{13,17,23,42,101\}$; per-seed std $\le 0.002$ on every metric); the
deterministic lexical, SapBERT, and BioSyn methods are seed-invariant, and the
train-seen ($^{\dagger}$) student rows are single-seed.}
\label{tab:results}
\end{table}

Three patterns hold across the table. On \texttt{alias\_name}, where surface
overlap is minimal, SapBERT dense lifts MRR from $0.042$ (lexical) to $0.316$
because the semantic encoder resolves descriptive names that share no characters
with the symbol. On \texttt{prev\_symbol}, where overlap is high, the lexical
baseline is already competitive ($0.205$) and the BioSyn hybrid is best ($0.250$),
confirming that sparse and dense signals are complementary when both are
informative. On the fact-disjoint \texttt{alias\_symbol} task every method is low
(MRR $\leq 0.134$): short alias symbols carry little semantic context and often no
shared substring, so neither a general biomedical embedding nor character overlap
resolves them reliably---this is where a notation-invariant encoder has the most
headroom.

\section{Analysis}
\label{sec:analysis}
The char-CNN student does not beat any reference method on the task it was added to
contest. It scores $0.059 \pm 0.001$ MRR on \texttt{alias\_symbol} (mean $\pm$ std
over $5$ seeds; best seed $0.060$), below SapBERT dense
($0.134$) and even the character $n$-gram floor ($0.077$)---a gap of more than $15$
seed standard deviations, so the ranking is not an artifact of a single training
run. Two factors explain the
failure. First, \texttt{alias\_symbol} is a held-out \emph{notation}, not merely
held-out facts: the student never sees an \texttt{alias\_symbol} form and must
transfer from other short-symbol notations whose surface statistics only partly
overlap. Second, a $\sim$100k-parameter byte-level char-CNN encodes little beyond
character composition; on short symbols that share few characters with the approved
symbol it has neither the lexical precision of TF--IDF $n$-grams nor the biomedical
semantics SapBERT acquired from UMLS self-alignment. The $0.985$ in-domain
validation top-1 confirms the model fit its training distribution well---it learned
a same-class surface similarity that does not transfer to an unseen short-symbol
notation against a $45$k-way pool.

The train-seen rows corroborate this reading. Even where the query notation was in
the training distribution, the small encoder trails the off-the-shelf baselines:
it carries some lexical signal on \texttt{prev\_symbol} (MRR $0.158$, which overlaps
approved symbols) but collapses to the lexical floor on free-text
\texttt{alias\_name} (MRR $0.042$), where it has no semantic capacity. The negative
result is therefore not an artifact of the held-out notation alone; the capacity of
a surface-only encoder is the deeper limit.

\subsection{Does attention distillation help?}
\label{sec:ablation}
Because the compression paradigm that motivates this work trains a student to
match a stronger model's \emph{internal} behavior---attention in particular
\citep{zagoruyko2017attention}---we test whether adding an attention-distillation
objective rescues the student. We compare three arms with identical student
configuration (char-CNN, hidden $64$, projection $128$, $3$ epochs,
same fact-disjoint split), each trained across $5$ random seeds
($\{13,17,23,42,101\}$), varying only the auxiliary loss that supervises the
student's cross-form character attention $\mathrm{softmax}(f_s f_t^{\top}/\sqrt{d})$:
(i) contrastive only; (ii) a \emph{deterministic} byte/digit teacher that places
mass on exact-character and digit--digit alignments; and (iii) a \emph{neural}
teacher---a frozen, higher-capacity char-CNN (hidden $128$, $6$ epochs, validation
top-1 $0.677$ versus the student's $0.582$) whose \emph{learned} attention is the
KL target. The distillation loss is active in both (ii) and (iii)
(train attention loss $>0$ throughout).

\begin{table}[t]
\centering
\small
\setlength{\tabcolsep}{4pt}
\begin{tabular}{lccc}
\toprule
\textbf{Student objective} & \textbf{MRR} & \textbf{top-1} & \textbf{inv.\ ratio} \\
\midrule
Contrastive only                 & $0.055 \pm 0.002$ & $0.036 \pm 0.003$ & $0.789 \pm 0.010$ \\
\;\; + byte-rule attn.\ distill   & $0.056 \pm 0.002$ & $0.037 \pm 0.004$ & $0.750 \pm 0.042$ \\
\;\; + neural-teacher attn.\ distill & $0.056 \pm 0.001$ & $0.038 \pm 0.002$ & $0.761 \pm 0.026$ \\
\bottomrule
\end{tabular}
\caption{Attention-distillation ablation on held-out \texttt{alias\_symbol}
(same $2{,}000$ queries / $45$k pool as Table~\ref{tab:results}), reported as
mean $\pm$ standard deviation across $5$ random seeds. The MRR difference between
the contrastive baseline and either distillation arm ($\sim\!0.001$) is well
within one standard deviation of seed variance. The invariance ratio is the
mean within-class over between-class embedding cosine distance (lower is more
notation-invariant); the three arms are indistinguishable once seed variance is
accounted for. Neither teacher improves held-out retrieval.}
\label{tab:ablation}
\end{table}

Neither form of attention distillation closes the gap (Table~\ref{tab:ablation}).
Across five seeds the deterministic and neural teachers both raise mean held-out
MRR by only $\sim\!0.001$ over the contrastive-only baseline ($0.055$)---a shift
smaller than the per-arm seed standard deviation, and therefore not distinguishable
from noise. The neural teacher carries a demonstrably stronger alignment signal
(validation top-1 $0.677$ versus the student's $0.582$) yet confers no measurable
retrieval benefit, and the three arms' invariance ratios overlap once seed variance
is included. Distilling a better \emph{surface-form} attention map does not induce
notation invariance here---consistent with the reading that the limit is
representational capacity, not the supervisory target's quality. This is a
test of attention-map matching between two surface-only encoders; it does not test
distillation from a \emph{semantically pretrained} teacher (e.g.\ SapBERT-strength
embeddings), which remains the open path (\S\ref{sec:conclusion}).

\section{Conclusion}
\label{sec:conclusion}
A lightweight contrastive char-CNN student does not, on its own, yield a
transferable cross-notation biomedical representation: it is beaten by both a
lexical floor and a general biomedical embedding on fact- and notation-disjoint
alias normalization. The result is only visible under a disjoint split; an
in-distribution measurement ($0.985$ top-1) would have suggested the opposite.
Closing the \texttt{alias\_symbol} gap requires what the current student
deliberately omits: a semantically pretrained backbone, distillation from a
SapBERT- or teacher-strength \emph{semantic} signal \citep{hinton2015distilling,
tian2020crd, hofstatter2020improving}---as opposed to the surface-form attention
teachers we ablate in \S\ref{sec:ablation}, which do not help---or exposure to the
target notation with notation-balanced sampling. We release the benchmark, the four reference methods,
and the student adapter to support that work.

\section*{Limitations}
The headline benchmark covers a single entity type (genes, via HGNC) with a
lightweight student; the numeric probe adds one further type (variants, via HGVS
cross-notation retrieval) but is likewise confined to a small char-CNN. Whether the
negative result extends to larger surface-only encoders or to other entity types
(chemicals, phenotypes) is not established here. In particular, the clinical
lab-value instantiation of the numeric half---held-out unit conversions and
reference-range normalization over real EHR values, which would require gated
PhysioNet/MIMIC-IV access---is left to future work; we make no lab-value claim.
Queries are capped at $2{,}000$ per notation for tractable CPU
evaluation, which bounds the precision of the reported metrics. The canonical
teacher is a structured oracle and is reported only to bracket the achievable
range, not as a learnable target.

Finally, the held-out \texttt{alias\_symbol} task is not perfectly well-posed from
surface form alone: some alias symbols are genuinely ambiguous or share no
substring with their approved symbol, so no notation-invariant encoder---however
large---can rank every query correctly from the string alone. The reference
methods bound how much of the task \emph{is} resolvable with lexical or dense
biomedical signal ($0.077$ and $0.134$ MRR respectively), and the low absolute
ceiling reflects this partial identifiability rather than any single model's
failure; our claim is only the relative one, that the char-CNN student does not
reach even that resolvable portion.

\section*{Ethics Statement}
All experiments use public reference terminologies (HGNC gene nomenclature) and
contain no patient data or personally identifiable information. The work is a
methods study on entity normalization and poses no direct dual-use concern.

\bibliography{references}

\end{document}
